\documentclass[11pt,a4paper]{article}

% ======================================================================
%  EagleEye — Cloud Architecture & Technology Stack (Technical Report)
% ======================================================================
\usepackage[T1]{fontenc}
\usepackage[utf8]{inputenc}
\usepackage{lmodern}
\usepackage[a4paper,margin=2.3cm,headheight=15pt]{geometry}
\usepackage{xcolor}
\usepackage{booktabs}
\usepackage{tabularx}
\usepackage{longtable}
\usepackage{enumitem}
\usepackage{listings}
\usepackage{fancyhdr}
\usepackage{titlesec}
\usepackage{parskip}
\usepackage[hidelinks]{hyperref}

% ---- palette ----
\definecolor{eagleblue}{HTML}{1F4E79}
\definecolor{eagleteal}{HTML}{2196F3}
\definecolor{eaglegray}{HTML}{555555}
\definecolor{codebg}{HTML}{F4F6F8}
\definecolor{rulegray}{HTML}{D0D5DB}
\definecolor{freegreen}{HTML}{2E7D32}
\definecolor{paidorange}{HTML}{E65100}

\hypersetup{pdftitle={EagleEye Cloud Architecture and Technology Stack},
  pdfauthor={EagleEye FYP}, colorlinks=true, linkcolor=eagleblue, urlcolor=eagleteal}

\titleformat{\section}{\Large\bfseries\color{eagleblue}}{\thesection}{0.6em}{}
\titleformat{\subsection}{\large\bfseries\color{eagleblue}}{\thesubsection}{0.6em}{}
\titlespacing*{\section}{0pt}{1.3em}{0.5em}

\pagestyle{fancy}
\fancyhf{}
\renewcommand{\headrulewidth}{0.4pt}
\renewcommand{\footrulewidth}{0.4pt}
\fancyhead[L]{\small\color{eaglegray}EagleEye --- Cloud Technology Stack}
\fancyhead[R]{\small\color{eaglegray}\leftmark}
\fancyfoot[L]{\small\color{eaglegray}Final Year Project --- Technical Report}
\fancyfoot[R]{\small\color{eaglegray}\thepage}

\lstdefinestyle{eagle}{backgroundcolor=\color{codebg},basicstyle=\ttfamily\small,
  breaklines=true,frame=single,rulecolor=\color{rulegray},framesep=6pt,xleftmargin=8pt,
  showstringspaces=false,columns=fullflexible,keepspaces=true}
\lstset{style=eagle}

\newcommand{\code}[1]{\texttt{#1}}
\newcommand{\free}{\textbf{\color{freegreen}FREE}}
\newcommand{\paid}{\textbf{\color{paidorange}PAID}}
\newcommand{\freepaid}{\textbf{\color{freegreen}FREE} tier \textbf{/} \textbf{\color{paidorange}PAID} upgrades}

% reusable technology header block
\newcommand{\techhead}[2]{\par\vspace{0.3em}\noindent
  \begin{tabularx}{\linewidth}{@{}lX@{}}
  \textbf{Type:} & #1 \\
  \textbf{Cost:} & #2 \\
  \end{tabularx}\par\vspace{0.3em}}

\begin{document}

% ====================== TITLE ======================
\begin{titlepage}
\centering
\vspace*{2cm}
{\color{eagleblue}\rule{\linewidth}{1.4pt}}\\[0.5cm]
{\Huge\bfseries\color{eagleblue} EagleEye}\\[0.35cm]
{\LARGE Cloud Architecture \& Technology Stack}\\[0.25cm]
{\large\color{eaglegray} A Detailed Report on the Technologies Used, Their Roles, and Their Costs}\\[0.4cm]
{\color{eagleblue}\rule{\linewidth}{1.4pt}}\\[1.2cm]

\begin{minipage}{0.86\linewidth}\centering\itshape\color{eaglegray}
EagleEye is a smart AI surveillance system built around an ESP32--CAM. This report
explains the migration from a local (LAN) architecture to a fully cloud-connected,
remote-capable architecture, and documents every technology in the new stack ---
what it does, why it was chosen, and whether it is free or paid.
\end{minipage}\\[1.6cm]

\begin{tabular}{rl}
\textbf{Project} & EagleEye --- AI Surveillance System \\[2pt]
\textbf{Core device} & ESP32--CAM (AI Thinker) \\[2pt]
\textbf{Document} & Cloud Technology Stack Report \\[2pt]
\textbf{Scope} & Architecture, technologies, pricing, comparison \\
\end{tabular}
\vfill
{\color{eaglegray}\small Pricing figures reflect publicly advertised free/paid tiers at the time of writing and should be re-verified before deployment.}
\end{titlepage}

\tableofcontents
\newpage

% ====================== 1. INTRODUCTION ======================
\section{Introduction}
EagleEye is an embedded AI surveillance node: an ESP32--CAM runs an on-device neural network
that detects humans, captures an image of an intruder, drives a servo to pan the camera, and
streams live video on request. The original prototype worked only on a \textbf{single local
Wi-Fi network} and depended on a PC. This report documents the redesign into a
\textbf{cloud-connected product} where the camera can be at one site and the phone anywhere in
the world.

\subsection{The architectural shift}
The single most important change is the \textbf{direction of communication}. In the old design
the phone had to reach \emph{into} the camera using its local IP address. In the new design the
camera connects \emph{outward} to cloud services and stays connected --- exactly how a messaging
app receives messages on a phone anywhere without anyone knowing the phone's IP. Because the
device initiates every connection, there is no need for a fixed IP, port forwarding, or a PC on
site, and the system works across any internet connection (home Wi-Fi, office, or 4G/LTE).

\subsection{What this report covers}
Section~\ref{sec:overview} gives a high-level before/after picture. Section~\ref{sec:stack}
lists the full stack at a glance. Section~\ref{sec:tech} documents \textbf{each technology} in
detail with its role and pricing. Section~\ref{sec:cost} summarises the total cost.
Section~\ref{sec:compare} compares old vs new, and Section~\ref{sec:conclusion} concludes.

% ====================== 2. OVERVIEW ======================
\section{Architecture Overview}\label{sec:overview}

\subsection{Old architecture (local / LAN)}
\begin{lstlisting}
   SITE (one local Wi-Fi)
 ESP32-CAM + Helper-ESP32(servo) --local MQTT--> PC (Mosquitto + bridge.py)
                                                 |--> Cloudinary (image)
                                                 |--> Firebase (record)
 Phone --- same Wi-Fi, types camera IP ---> ws://<ip>:81 (video), http://<ip>/servo
\end{lstlisting}
\textbf{Problems:} works only on the same Wi-Fi; the camera IP must be typed and changes
often; a PC must stay powered to run the broker and the upload bridge; a second ESP32 board is
needed just for the servo; the local broker is unencrypted and open.

\subsection{New architecture (cloud)}
\begin{lstlisting}
   SITE A                       CLOUD                         ANYWHERE
 ESP32-CAM (AI + servo) --TLS--> HiveMQ broker <--wss-- Phone app (control)
        |  image  --HTTPS--> Cloudinary
        |  alert  --REST---> Firebase  ----------------> Phone app (history)
        |  video  --wss----> Deno relay ---------------> Phone app (live, on-demand)
\end{lstlisting}
\textbf{Result:} no PC, no helper board, no LAN IP. The camera reaches out to the cloud;
the phone reaches the same cloud; they meet in the middle. It works from any network and is
encrypted end-to-end.

% ====================== 3. STACK AT A GLANCE ======================
\section{The Technology Stack at a Glance}\label{sec:stack}
\renewcommand{\arraystretch}{1.25}
\begin{tabularx}{\linewidth}{@{}lXl@{}}
\toprule
\textbf{Technology} & \textbf{Role in EagleEye} & \textbf{Cost} \\
\midrule
MQTT (protocol)        & Lightweight messaging for commands/status/alerts & \free \\
HiveMQ Cloud           & Managed MQTT broker (the cloud "post office")     & \freepaid \\
TLS / Let's Encrypt    & Encryption + certificates for all connections     & \free \\
Cloudinary             & Stores intruder images, returns a URL             & \freepaid \\
Firebase Realtime DB   & Stores the alert history the app reads            & \freepaid \\
Deno \& Deno Deploy    & Hosts the live-video relay server                 & \freepaid \\
WebSocket (protocol)   & Transport for relayed video + app$\leftrightarrow$broker & \free \\
React Native + Expo    & The cross-platform mobile app                     & \freepaid \\
mqtt.js                & MQTT client library inside the app                & \free \\
ESP32-CAM + Arduino    & The camera hardware + firmware framework          & hardware only \\
Edge Impulse           & Trains / exports the human-detection AI model     & \freepaid \\
ESP32 Arduino libs     & PubSubClient, ArduinoJson, WebSockets, ESP32Servo & \free \\
\bottomrule
\end{tabularx}

% ====================== 4. TECHNOLOGIES IN DETAIL ======================
\section{Technologies in Detail}\label{sec:tech}

% ---------- MQTT ----------
\subsection{MQTT (Message Queuing Telemetry Transport)}
\techhead{Open messaging protocol / standard}{\free (the protocol is royalty-free)}
\textbf{What it is.} MQTT is a lightweight publish/subscribe messaging protocol designed for
IoT devices with limited bandwidth and power. Devices connect to a central \emph{broker} and
exchange small messages on named \emph{topics} (mailboxes). A device can \emph{publish} to a
topic and/or \emph{subscribe} to receive everything posted there.

\textbf{Role in EagleEye.} MQTT is the control backbone. The camera publishes its
\code{status} and \code{alert} topics and subscribes to a \code{cmd} topic; the app publishes
commands (arm/disarm, servo angle, start/stop stream). Two features are especially useful:
\emph{retained messages} (the broker keeps the last status so the app sees current state
instantly) and the \emph{Last Will and Testament} (the broker auto-announces the camera is
offline if it disconnects unexpectedly).

\textbf{Free or paid?} The protocol itself is an open OASIS/ISO standard --- completely
\free. Costs only arise from the broker that hosts it (next subsection).

\textbf{Why chosen.} It is the de-facto industry standard for IoT command-and-control (used by
AWS IoT, Azure IoT Hub, etc.), extremely light on the ESP32, and naturally outbound-connecting.

% ---------- HiveMQ ----------
\subsection{HiveMQ Cloud}
\techhead{Managed cloud MQTT broker (SaaS)}{\freepaid}
\textbf{What it is.} HiveMQ Cloud is a hosted, always-on MQTT broker. Instead of running our
own broker on a PC, we use HiveMQ's managed service, which provides a secure endpoint, user
credentials, and per-topic access control.

\textbf{Role in EagleEye.} It is the cloud "post office" both the camera and the app connect
to. The camera connects over TLS (port 8883); the app connects over secure WebSocket (port
8884). All control, status, and alert signalling flows through it.

\textbf{Free or paid?} HiveMQ Cloud offers a \free \textbf{Serverless} tier (sufficient for
this project --- it supports a small number of devices and a monthly data allowance at no cost
and requires no credit card). \paid plans (Starter / Professional / dedicated clusters) scale
to more connections, more data throughput, and enterprise features. EagleEye runs entirely on
the free tier.

\textbf{Why chosen.} Zero-maintenance, secure by default (TLS + per-credential access), free
for our scale, and removes the PC-hosted broker entirely.

% ---------- TLS ----------
\subsection{TLS / Let's Encrypt}
\techhead{Transport-layer encryption + certificate authority}{\free}
\textbf{What it is.} TLS (Transport Layer Security) is the same encryption that powers
\code{https://}. It encrypts data in transit so no one on the network can read or tamper with
it. Certificates (often issued free by Let's Encrypt) let the client verify it is talking to
the genuine server.

\textbf{Role in EagleEye.} Every cloud connection is encrypted: MQTT over TLS to HiveMQ,
HTTPS to Cloudinary and Firebase, and secure WebSocket (\code{wss}) to the relay. This replaces
the old open, unencrypted local broker.

\textbf{Free or paid?} \free --- TLS is built into the libraries, and certificates from
Let's Encrypt (used by the cloud providers) are free.

\textbf{Why chosen.} Security is mandatory once data leaves the local network; TLS is the
universal standard and is provided automatically by every service we use.

% ---------- Cloudinary ----------
\subsection{Cloudinary}
\techhead{Cloud image storage \& delivery (media CDN)}{\freepaid}
\textbf{What it is.} Cloudinary is a media-management cloud that stores images and serves them
over a fast global URL. It supports \emph{unsigned upload presets}, which let a device upload
directly without holding a secret key.

\textbf{Role in EagleEye.} When the AI detects an intruder, the camera uploads the captured
JPEG \emph{directly} to Cloudinary over HTTPS and receives back a public image URL. That URL is
then stored in the alert record. This removed the PC bridge that used to do the upload.

\textbf{Free or paid?} Cloudinary has a \free plan with a monthly allowance (storage,
bandwidth, and transformations measured in "credits"), which is ample for a surveillance demo.
\paid plans add more storage/bandwidth and advanced features. EagleEye uses the free plan with
an unsigned upload preset scoped to one folder.

\textbf{Why chosen.} Free, gives an instant shareable URL the app can display, and supports
secure device-side uploads without embedding secret keys in the firmware.

% ---------- Firebase ----------
\subsection{Firebase Realtime Database}
\techhead{Cloud NoSQL real-time database (Google)}{\freepaid}
\textbf{What it is.} Firebase Realtime Database (RTDB) is a cloud-hosted JSON database that
pushes updates to connected clients in real time. When data changes, every listening app
updates instantly.

\textbf{Role in EagleEye.} It stores the \textbf{alert history}. The camera writes each alert
(\code{timestamp}, \code{image\_url}, \code{type}) into the \code{alerts} node via the REST
API; the mobile app listens to that node and shows new intrusions immediately in the dashboard
and gallery.

\textbf{Free or paid?} The \free \textbf{Spark} plan covers our needs (a generous storage and
download allowance and a limited number of simultaneous connections). The \paid
\textbf{Blaze} (pay-as-you-go) plan is only needed for higher volume or for Cloud Functions /
push notifications. EagleEye's alert history runs on the free tier.

\textbf{Why chosen.} Real-time updates with no server code, an easy REST API the ESP32 can call
directly, and it was already integrated in the original app, so the alert UI needed no changes.

% ---------- Deno ----------
\subsection{Deno \& Deno Deploy}
\techhead{JavaScript/TypeScript runtime + serverless edge hosting}{\freepaid}
\textbf{What it is.} Deno is a modern, secure runtime for JavaScript/TypeScript (an alternative
to Node.js). \emph{Deno Deploy} is its serverless hosting platform that runs code on edge
servers worldwide with built-in HTTPS/WebSocket support.

\textbf{Role in EagleEye.} It hosts the \textbf{live-video relay} --- a small program that acts
as a "meeting room": the camera connects out and pushes JPEG frames; the phone connects out and
receives them. The relay simply forwards frames, so neither side needs the other's IP. This is
what makes remote live video possible.

\textbf{Free or paid?} Deno Deploy has a \free tier with a monthly request and bandwidth
allowance that comfortably covers on-demand surveillance streaming. \paid (Pro) plans raise
those limits. The relay runs on the free tier.

\textbf{Why chosen.} Free, always-on, supports WebSockets natively, deploys a single file in
minutes, and terminates TLS automatically (giving a \code{wss://} URL with no certificate work).

% ---------- WebSocket ----------
\subsection{WebSocket Protocol}
\techhead{Full-duplex persistent network protocol / standard}{\free}
\textbf{What it is.} WebSocket is a protocol that keeps a single connection open for continuous
two-way data, unlike normal HTTP request/response. It is ideal for streaming and live updates.

\textbf{Role in EagleEye.} It carries the \textbf{live video frames} between the camera, the
relay, and the app, and also carries MQTT messages between the app and HiveMQ
(MQTT-over-WebSocket). Persistent connections mean low latency and no constant polling.

\textbf{Free or paid?} \free --- an open web standard built into browsers, React Native, and
the ESP32 libraries.

\textbf{Why chosen.} It is the natural transport for continuous binary frames and is supported
everywhere (device, relay, and phone), letting the same decode/render path be reused.

% ---------- React Native + Expo ----------
\subsection{React Native \& Expo}
\techhead{Cross-platform mobile app framework + tooling}{\freepaid}
\textbf{What it is.} React Native is an open-source framework (by Meta) for building native
Android and iOS apps from a single JavaScript codebase. Expo is a layer on top that simplifies
development, testing (via the Expo Go app), and building.

\textbf{Role in EagleEye.} It is the entire \textbf{mobile app}: the dashboard (arm/disarm,
online status, latest alert), the live-view screen with the servo pan slider, the gallery, and
the statistics charts. One codebase runs on both Android and iOS.

\textbf{Free or paid?} React Native and the Expo SDK are open-source and \free. Testing in
\textbf{Expo Go} is free. Expo's \textbf{EAS} cloud build/submit/push service has a \free tier
(a limited number of builds per month) with \paid plans for higher build volume and priority.
EagleEye is developed and demoed for free in Expo Go.

\textbf{Why chosen.} One codebase for both platforms, fast iteration with live reload, large
ecosystem, and a free path from development to a shippable build.

% ---------- mqtt.js ----------
\subsection{mqtt.js}
\techhead{Open-source MQTT client library (JavaScript)}{\free}
\textbf{What it is.} \code{mqtt.js} is the most widely used JavaScript MQTT client. It runs in
browsers and React Native and speaks MQTT over WebSocket.

\textbf{Role in EagleEye.} It is the app's connection to HiveMQ: it publishes commands
(arm, servo, stream) and subscribes to the camera's status and stream-ready signals. It
replaced the old "type the camera's IP" approach for control.

\textbf{Free or paid?} \free --- open-source (MIT/EPL).

\textbf{Why chosen.} Pure JavaScript (works in Expo Go without native modules), maintained, and
the standard choice for MQTT in web/mobile apps.

% ---------- ESP32-CAM ----------
\subsection{ESP32-CAM \& the Arduino Framework}
\techhead{Microcontroller-with-camera hardware + open firmware framework}{Hardware cost only; software \free}
\textbf{What it is.} The ESP32--CAM (AI Thinker) is a low-cost board combining a dual-core
ESP32 (Wi-Fi + processing) with an OV2640 camera and PSRAM. It is programmed with the
open-source ESP32 Arduino core.

\textbf{Role in EagleEye.} It is the \textbf{device itself} --- it runs the AI, captures
images, drives the servo on GPIO~15, connects to the cloud, and streams video. In the new
design it also took over the servo job (no second board).

\textbf{Free or paid?} The \textbf{hardware} is a one-time purchase (roughly \$6--\$10 per
board, plus a USB programmer and a 5\,V supply). The \textbf{software} --- the Arduino core and
all toolchains --- is \free and open-source.

\textbf{Why chosen.} Extremely low cost for a Wi-Fi camera with enough power to run a small
neural network on-device, with a huge community and free tooling.

% ---------- Edge Impulse ----------
\subsection{Edge Impulse (On-device AI)}
\techhead{TinyML model training \& deployment platform}{\freepaid}
\textbf{What it is.} Edge Impulse is a platform for collecting data, training small neural
networks, and exporting them as optimised C++ libraries that run on microcontrollers (TinyML).

\textbf{Role in EagleEye.} It produced the \textbf{human-detection model} (an RGB
$96\times96$ classifier) that runs \emph{on the ESP32 itself} --- no cloud inference, no
per-image AI cost. The exported library performs the resize, normalisation, and neural-network
inference on-device every frame.

\textbf{Free or paid?} Edge Impulse has a \free \textbf{Developer/Community} plan that covers
training and exporting models for personal/educational projects. \paid \textbf{Enterprise}
plans add collaboration, larger datasets, and commercial features. The exported model runs
on-device for free forever. EagleEye uses the free plan.

\textbf{Why chosen.} It makes on-device AI accessible without deep ML expertise, and on-device
inference means privacy (images analysed locally) and zero ongoing AI cost.

% ---------- Arduino libs ----------
\subsection{Supporting ESP32 Arduino Libraries}
\techhead{Open-source firmware libraries}{\free}
These libraries (all \free, open-source) implement the device side of each technology:
\begin{itemize}[nosep,leftmargin=1.4em]
  \item \textbf{PubSubClient} --- the MQTT client on the ESP32 (publishes status/alerts,
        receives commands), used over a TLS connection.
  \item \textbf{ArduinoJson} --- builds and parses the JSON messages for commands, status, and
        alert records.
  \item \textbf{WebSockets (by Markus Sattler)} --- the WebSocket client that pushes video
        frames to the Deno relay.
  \item \textbf{ESP32Servo} --- generates the servo control signal on GPIO~15, using a timer
        that does not conflict with the camera.
  \item \textbf{esp32-camera} --- the camera driver (capture + JPEG conversion).
\end{itemize}
\textbf{Why chosen.} They are the standard, well-maintained, free building blocks of the ESP32
Arduino ecosystem and map one-to-one onto the cloud services above.

% ====================== 5. COST ANALYSIS ======================
\section{Cost Analysis: Free vs Paid}\label{sec:cost}
A central design goal was that the entire cloud stack runs on \textbf{free tiers}, so the only
real cost is the hardware. The table below summarises each item.

\renewcommand{\arraystretch}{1.3}
\begin{tabularx}{\linewidth}{@{}lXl@{}}
\toprule
\textbf{Item} & \textbf{Free tier used by EagleEye} & \textbf{When you would pay} \\
\midrule
HiveMQ Cloud      & Serverless free tier (small fleet, monthly data) & Many devices / high message volume \\
Cloudinary        & Free plan (monthly credit allowance)             & Large media storage / bandwidth \\
Firebase RTDB     & Spark free plan                                  & High traffic, Cloud Functions, push \\
Deno Deploy       & Free tier (requests + bandwidth)                 & Heavy / continuous streaming \\
Expo (EAS)        & Expo Go + free build allowance                   & Frequent production builds \\
Edge Impulse      & Developer/Community plan                         & Commercial / enterprise use \\
MQTT, TLS, WebSocket, RN, mqtt.js, Arduino libs & Open standards \& open-source & Never (always free) \\
\midrule
\textbf{Hardware} & ESP32--CAM + programmer + 5\,V supply + servo    & One-time purchase ($\approx$\$15--\$25 total) \\
\bottomrule
\end{tabularx}

\vspace{0.5em}
\noindent\textbf{Bottom line.} For a final-year-project scale (one camera, one or two users),
the \textbf{recurring software cost is \$0/month}. The only spend is the one-time hardware. The
architecture scales to paid tiers smoothly if it were ever commercialised, without changing the
design.

% ====================== 6. OLD VS NEW ======================
\section{Old vs New: Summary of Changes}\label{sec:compare}
\renewcommand{\arraystretch}{1.3}
\begin{tabularx}{\linewidth}{@{}lXX@{}}
\toprule
\textbf{Aspect} & \textbf{Old (LAN)} & \textbf{New (Cloud)} \\
\midrule
Reach           & Same Wi-Fi only            & Anywhere via the internet \\
Addressing      & Type the camera's IP       & Device ID; no IP needed \\
Broker          & Mosquitto on a PC, no TLS  & HiveMQ Cloud, TLS-secured \\
Image upload    & PC bridge (\code{bridge.py}) & Camera uploads directly to Cloudinary \\
Alert storage   & Written by the PC bridge   & Camera writes Firebase directly \\
Live video      & Direct to camera IP        & Cloud relay, on-demand \\
Control         & Firebase + bridge relay    & Direct MQTT commands \\
Online status   & Heartbeat guesswork        & MQTT retained status + Last Will \\
Servo           & Separate "helper" ESP32    & Driven on the main board (GPIO~15) \\
PC required     & Yes                        & No \\
Security        & Open, unencrypted          & Encrypted + per-device credentials \\
\bottomrule
\end{tabularx}

% ====================== 7. CONCLUSION ======================
\section{Conclusion}\label{sec:conclusion}
The redesign of EagleEye replaced a fragile, location-locked, PC-dependent system with a
clean, secure, cloud-connected one. The guiding principle --- \emph{the device reaches out to
the cloud instead of waiting to be reached} --- eliminated the manual-IP problem, the on-site
PC, and the extra microcontroller, while adding encryption, a truthful online/offline
indicator, push-style alerting, and the ability to operate from anywhere.

Equally important for a student project, every cloud technology in the stack was selected to
run on a \textbf{free tier}: MQTT, TLS, and WebSocket are open standards; HiveMQ, Cloudinary,
Firebase, Deno Deploy, Expo, and Edge Impulse all provide capable free plans; and the firmware
libraries are open-source. The result is a professional, industry-aligned architecture
(mirroring how commercial IoT platforms such as AWS IoT and Ring/Nest cameras are built) at
\textbf{zero recurring software cost}, with a clear, unchanged path to paid tiers should the
system ever be scaled to many cameras and users.

\vspace{1em}
\noindent{\color{rulegray}\rule{\linewidth}{0.4pt}}\\
{\small\itshape\color{eaglegray} EagleEye FYP --- Cloud Architecture \& Technology Stack report.
Free/paid tiers are based on publicly advertised pricing at the time of writing; verify current
limits with each provider before relying on them.}

\end{document}
